% File: ai-nhap/2026-11-15_17_de_tham_khao_go_vap_de_1_va_go_vap_de_2.tex
% Nguồn: Tuyển tập đề tham khảo môn Toán Tuyển sinh 10 TP.HCM 2025-2026
% Gồm 14 bài: Quận Gò Vấp - Đề 1 (7 bài) và Quận Gò Vấp - Đề 2 (7 bài)

% ==============================================================================
% PHẦN 1: ĐỀ THAM KHẢO TUYỂN SINH 10 QUẬN GÒ VẤP – ĐỀ 1 (7 BÀI)
% ==============================================================================

\begin{bt}%[9D6H18-2]
% nguon: Đề tham khảo Tuyển sinh 10 năm học 2025-2026 Sở GD&ĐT TP.HCM – Quận Gò Vấp - Đề 1 (Bài 1)
Cho parabol $(P)\colon y = -2x^2$.
\begin{enumerate}[a)]
    \item Vẽ đồ thị $(P)$ trên hệ trục tọa độ $Oxy$.
    \item Tìm tọa độ các điểm $M$ thuộc $(P)$ có hoành độ bằng tung độ.
\end{enumerate}
\loigiai{
\begin{enumerate}[a)]
    \item Bảng giá trị tương ứng của hàm số $y = -2x^2$:
    \begin{center}
    \begin{tabular}{|c|c|c|c|c|c|}
        \hline
        $x$ & $-2$ & $-1$ & $0$ & $1$ & $2$ \\
        \hline
        $y = -2x^2$ & $-8$ & $-2$ & $0$ & $-2$ & $-8$ \\
        \hline
    \end{tabular}
    \end{center}
    Đồ thị $(P)$ là một parabol có đỉnh $O(0; 0)$, đối xứng qua trục tung $Oy$ và bề lõm quay xuống dưới.
    \item Điểm $M(x; y)$ thuộc $(P)$ có hoành độ bằng tung độ khi và chỉ khi:
    $x = y \iff -2x^2 = x \iff 2x^2 + x = 0 \iff x(2x + 1) = 0 \iff \left[\begin{aligned} &x = 0 \implies y = 0 \\ &x = -\dfrac{1}{2} \implies y = -\dfrac{1}{2}. \end{aligned}\right.$
    Vậy có hai điểm thỏa mãn yêu cầu đề bài là $M_1(0; 0)$ và $M_2\left(-\dfrac{1}{2}; -\dfrac{1}{2}\right)$.
\end{enumerate}
}
\end{bt}

\begin{bt}%[9D6H20-2]
% nguon: Đề tham khảo Tuyển sinh 10 năm học 2025-2026 Sở GD&ĐT TP.HCM – Quận Gò Vấp - Đề 1 (Bài 2)
Cho phương trình bậc hai $x^2 - 5x + 6 = 0$.
\begin{enumerate}[a)]
    \item Chứng minh phương trình luôn có hai nghiệm phân biệt $x_1, x_2$.
    \item Không giải phương trình, hãy tính giá trị của biểu thức $A = x_1 + \dfrac{x_2^2}{5}$.
\end{enumerate}
\loigiai{
\begin{enumerate}[a)]
    \item Phương trình $x^2 - 5x + 6 = 0$ có các hệ số $a = 1$, $b = -5$, $c = 6$.
    Biệt thức: $\Delta = (-5)^2 - 4 \cdot 1 \cdot 6 = 25 - 24 = 1 > 0$.
    Do đó phương trình luôn có hai nghiệm phân biệt $x_1, x_2$.
    \item Theo định lí Viète, ta có:
    $\begin{cases} S = x_1 + x_2 = 5 \\ P = x_1 x_2 = 6. \end{cases}$
    Vì $x_2$ là nghiệm của phương trình nên $x_2^2 - 5x_2 + 6 = 0 \implies x_2^2 = 5x_2 - 6$.
    Thay vào biểu thức $A$, ta được:
    $A = x_1 + \dfrac{5x_2 - 6}{5} = x_1 + x_2 - \dfrac{6}{5} = S - \dfrac{6}{5} = 5 - \dfrac{6}{5} = \dfrac{19}{5} = 3{,}8$.
\end{enumerate}
}
\end{bt}

\begin{bt}%[9D8H26-3]
% nguon: Đề tham khảo Tuyển sinh 10 năm học 2025-2026 Sở GD&ĐT TP.HCM – Quận Gò Vấp - Đề 1 (Bài 3)
Bác Dũng có một chiếc ổ khóa vali gồm $4$ vòng số quay độc lập, mỗi vòng gồm các chữ số từ $0$ đến $9$. Bác chọn ngẫu nhiên một dãy gồm $4$ chữ số để thiết lập làm mật mã mở khóa. Tính xác suất của các biến cố sau:
\begin{enumerate}[a)]
    \item Biến cố $A$: ``Bốn chữ số được chọn hoàn toàn giống nhau''.
    \item Biến cố $B$: ``Bốn chữ số được chọn lập thành một số tự nhiên có bốn chữ số''.
    \item Biến cố $C$: ``Bốn chữ số được chọn có tổng các chữ số bằng $35$''.
\end{enumerate}
\loigiai{
Mỗi vòng quay có $10$ chữ số (từ $0$ đến $9$). Số mật mã gồm $4$ chữ số có thể tạo thành là:
$n(\Omega) = 10^4 = 10\,000$ (kết quả đồng khả năng).
\begin{enumerate}[a)]
    \item Các mật mã có $4$ chữ số giống nhau là $\{0000; 1111; 2222; \dots; 9999\}$ ($10$ mật mã).
    Xác suất của biến cố $A$ là:
    $P(A) = \dfrac{n(A)}{n(\Omega)} = \dfrac{10}{10\,000} = \dfrac{1}{1000} = 0{,}001$.
    \item Mật mã lập thành một số tự nhiên có bốn chữ số khi chữ số đầu tiên khác $0$.
    Chữ số đầu tiên có $9$ cách chọn (từ $1$ đến $9$), ba chữ số sau mỗi chữ số có $10$ cách chọn.
    Số kết quả thuận lợi cho $B$ là: $n(B) = 9 \cdot 10^3 = 9000$.
    Xác suất của biến cố $B$ là:
    $P(B) = \dfrac{n(B)}{n(\Omega)} = \dfrac{9000}{10\,000} = \dfrac{9}{10} = 0{,}9$.
    \item Vì mỗi chữ số không vượt quá $9$ nên tổng lớn nhất của $4$ chữ số là $4 \cdot 9 = 36$.
    Để tổng bốn chữ số bằng $35$, bộ bốn chữ số bắt buộc phải gồm ba chữ số $9$ và một chữ số $8$ (vì $9 + 9 + 9 + 8 = 35$).
    Số hoán vị của bộ số $(9; 9; 9; 8)$ là $4$ mật mã: $8999, 9899, 9989, 9998$.
    Do đó $n(C) = 4$.
    Xác suất của biến cố $C$ là:
    $P(C) = \dfrac{n(C)}{n(\Omega)} = \dfrac{4}{10\,000} = \dfrac{1}{2500} = 0{,}0004$.
\end{enumerate}
}
\end{bt}

\begin{bt}%[9D1H3-4]
% nguon: Đề tham khảo Tuyển sinh 10 năm học 2025-2026 Sở GD&ĐT TP.HCM – Quận Gò Vấp - Đề 1 (Bài 4)
Trong cuộc thi ``Đố vui học tập'', mỗi học sinh phải trả lời $12$ câu hỏi trắc nghiệm. Mỗi câu trả lời đúng được cộng $5$ điểm, trả lời sai bị trừ $2$ điểm. Khi bước vào cuộc thi, mỗi thí sinh được ban tổ chức tặng sẵn $10$ điểm khởi đầu. Thí sinh nào đạt tổng kết quả từ $40$ điểm trở lên sẽ được quyền vào vòng thi tiếp theo. Gọi $x$ là số câu trả lời đúng của một thí sinh ($x \in \mathbb{N}, x \le 12$).
\begin{enumerate}[a)]
    \item Hãy thiết lập bất phương trình theo ẩn $x$ biểu diễn điều kiện để thí sinh đạt từ $40$ điểm trở lên.
    \item Hỏi thí sinh đó phải trả lời đúng ít nhất bao nhiêu câu hỏi để được bước tiếp vào vòng sau?
\end{enumerate}
\loigiai{
\begin{enumerate}[a)]
    \item Nếu thí sinh trả lời đúng $x$ câu thì số câu trả lời sai là $12 - x$ câu.
    Tổng số điểm thí sinh đạt được sau $12$ câu hỏi là:
    $T(x) = 10 + 5x - 2(12 - x) = 10 + 5x - 24 + 2x = 7x - 14$ (điểm).
    Bất phương trình biểu thị điều kiện để thí sinh đạt từ $40$ điểm trở lên là:
    $7x - 14 \ge 40$.
    \item Giải bất phương trình:
    $7x - 14 \ge 40 \iff 7x \ge 54 \iff x \ge \dfrac{54}{7} \approx 7{,}71$.
    Vì số câu trả lời đúng $x$ phải là số tự nhiên không vượt quá $12$ nên giá trị nhỏ nhất của $x$ là $8$.
    Vậy thí sinh phải trả lời đúng ít nhất $8$ câu hỏi thì mới được vào vòng thi tiếp theo.
\end{enumerate}
}
\end{bt}

\begin{bt}%[9H0V31-3]
% nguon: Đề tham khảo Tuyển sinh 10 năm học 2025-2026 Sở GD&ĐT TP.HCM – Quận Gò Vấp - Đề 1 (Bài 5)
Một công ty nước giải khát nghiên cứu sản xuất hai mẫu lon nhôm hình trụ: mẫu lon phổ biến có chiều cao $12\text{ cm}$ và đường kính đáy $6{,}5\text{ cm}$; mẫu lon thon cao có chiều cao $14\text{ cm}$ và đường kính đáy $6\text{ cm}$. Lấy $\pi \approx 3{,}1416$.
\begin{enumerate}[a)]
    \item Tính thể tích của mẫu lon thon cao và cho biết mẫu lon này có chứa được trọn vẹn $335\text{ ml}$ chất lỏng của lon phổ biến hay không (kết quả làm tròn đến hàng phần mười của $\text{cm}^3$)?
    \item Biết chi phí vật liệu sản xuất lon tỉ lệ thuận với diện tích toàn phần của lon. Hỏi chi phí sản xuất một chiếc lon thon cao tăng thêm bao nhiêu phần trăm so với chiếc lon phổ biến (kết quả làm tròn đến hàng phần mười)?
\end{enumerate}
\loigiai{
\begin{enumerate}[a)]
    \item Mẫu lon thon cao có bán kính đáy $r_2 = \dfrac{6}{2} = 3\text{ cm}$ và chiều cao $h_2 = 14\text{ cm}$.
    Thể tích của mẫu lon thon cao là:
    $V_2 = \pi r_2^2 h_2 = \pi \cdot 3^2 \cdot 14 = 126\pi \approx 395{,}8\text{ cm}^3 = 395{,}8\text{ ml}$.
    Vì $395{,}8\text{ ml} > 335\text{ ml}$ nên mẫu lon thon cao hoàn toàn chứa được hết lượng nước ngọt $335\text{ ml}$ của mẫu lon phổ biến.
    \item Mẫu lon phổ biến có bán kính đáy $r_1 = \dfrac{6{,}5}{2} = 3{,}25\text{ cm}$ và chiều cao $h_1 = 12\text{ cm}$.
    Diện tích toàn phần của mẫu lon phổ biến là:
    $S_{\text{tp1}} = 2\pi r_1 h_1 + 2\pi r_1^2 = 2\pi \cdot 3{,}25 \cdot 12 + 2\pi \cdot (3{,}25)^2 = 78\pi + 21{,}125\pi = 99{,}125\pi\text{ (cm}^2\text{)}$.
    Diện tích toàn phần của mẫu lon thon cao là:
    $S_{\text{tp2}} = 2\pi r_2 h_2 + 2\pi r_2^2 = 2\pi \cdot 3 \cdot 14 + 2\pi \cdot 3^2 = 84\pi + 18\pi = 102\pi\text{ (cm}^2\text{)}$.
    Tỉ lệ phần trăm chi phí sản xuất lon thon cao tăng thêm so với lon phổ biến là:
    $\dfrac{S_{\text{tp2}} - S_{\text{tp1}}}{S_{\text{tp1}}} \cdot 100\% = \dfrac{102\pi - 99{,}125\pi}{99{,}125\pi} \cdot 100\% = \dfrac{2{,}875}{99{,}125} \cdot 100\% \approx 2{,}9\%$.
\end{enumerate}
}
\end{bt}

\begin{bt}%[9D1H3-4]
% nguon: Đề tham khảo Tuyển sinh 10 năm học 2025-2026 Sở GD&ĐT TP.HCM – Quận Gò Vấp - Đề 1 (Bài 6)
Mẹ bạn An dự tính rằng sau $3$ năm nữa cần có số tiền $250\,000\,000$ đồng để chu cấp cho An theo học $4$ năm đại học. Mẹ quyết định đem gửi một khoản tiền tiết kiệm ban đầu vào ngân hàng theo hình thức lãi kép (tiền lãi hàng năm được tự động cộng dồn vào vốn để tính lãi cho năm tiếp theo) với lãi suất ổn định là $6{,}3\%$/năm. Hỏi ngay từ bây giờ mẹ bạn An cần gửi vào ngân hàng số tiền vốn ít nhất là bao nhiêu (kết quả làm tròn đến hàng chục nghìn đồng)?
\loigiai{
Công thức tính số tiền tích lũy cả gốc và lãi sau $n$ năm gửi tiết kiệm theo hình thức lãi kép là:
$A = P(1 + r)^n$.
Trong đó $A = 250\,000\,000$ đồng, lãi suất $r = 6{,}3\% = 0{,}063$ và kỳ hạn $n = 3$ năm.
Số tiền vốn gốc $P$ mẹ bạn An cần gửi ngay từ bây giờ là:
$P = \dfrac{A}{(1 + r)^n} = \dfrac{250\,000\,000}{(1 + 0{,}063)^3} = \dfrac{250\,000\,000}{(1{,}063)^3} \approx 208\,129\,660$ (đồng).
Làm tròn đến hàng chục nghìn đồng, số tiền mẹ bạn An cần gửi là $208\,130\,000$ đồng.
}
\end{bt}

\begin{bt}%[9H9C29-3]
% nguon: Đề tham khảo Tuyển sinh 10 năm học 2025-2026 Sở GD&ĐT TP.HCM – Quận Gò Vấp - Đề 1 (Bài 7)
Cho tam giác nhọn $ABC$ nội tiếp đường tròn $(O; R)$. Đường tròn tâm $I$ đường kính $BC$ cắt hai cạnh $AB$ và $AC$ lần lượt tại $E$ và $D$ ($E \ne B, D \ne C$). Hai đoạn thẳng $BD$ và $CE$ cắt nhau tại trực tâm $H$.
\begin{enumerate}[a)]
    \item Chứng minh $AH \perp BC$.
    \item Kẻ đường kính $AM$ của đường tròn $(O)$. Chứng minh ba điểm $H, I, M$ thẳng hàng.
    \item Trong trường hợp $\dfrac{ED}{BC} = \dfrac{1}{2}$, hãy tính độ dài đoạn thẳng $AH$ theo bán kính $R$.
\end{enumerate}
\loigiai{
\begin{enumerate}[a)]
    \item Vì $D, E$ thuộc đường tròn đường kính $BC$ nên $\widehat{BDC} = 90^\circ$ và $\widehat{BEC} = 90^\circ$, suy ra $BD \perp AC$ và $CE \perp AB$.
    Do đó $BD$ và $CE$ là hai đường cao của tam giác $ABC$.
    Vì $H$ là giao điểm của $BD$ và $CE$ nên $H$ là trực tâm của tam giác $ABC$, suy ra $AH \perp BC$.
    \item Vì $AM$ là đường kính của $(O)$ nên $\widehat{ABM} = 90^\circ$ và $\widehat{ACM} = 90^\circ$, suy ra $BM \perp AB$ và $CM \perp AC$.
    Lại có $CH \perp AB$ và $BH \perp AC$, do đó $CH \parallel BM$ và $BH \parallel CM$.
    Tứ giác $BHCM$ có các cặp cạnh đối song song nên là hình bình hành.
    Vì $I$ là trung điểm của đường chéo $BC$ nên $I$ cũng là trung điểm của đường chéo $HM$.
    Do đó ba điểm $H, I, M$ thẳng hàng (đpcm).
    \item Tứ giác $BEDC$ nội tiếp đường tròn đường kính $BC$ nên $\triangle ADE \backsim \triangle ABC$ (c.g.c).
    Tỉ số đồng dạng là:
    $\dfrac{ED}{BC} = \cos \widehat{BAC} = \dfrac{1}{2} \implies \widehat{BAC} = 60^\circ$.
    Vì tam giác $ABC$ có $\widehat{BAC} = 60^\circ$ nội tiếp $(O; R)$ nên góc ở tâm $\widehat{BOC} = 2\widehat{BAC} = 120^\circ$.
    Trong tam giác cân $OBC$ có $OB = OC = R$ và $\widehat{BOC} = 120^\circ$, ta có khoảng cách từ $O$ đến dây $BC$ là:
    $OI = R \cdot \cos 60^\circ = \dfrac{R}{2}$.
    Trong hình bình hành $BHCM$, đoạn $OI$ là đường trung bình của tam giác $AHM$, suy ra:
    $AH = 2OI = 2 \cdot \dfrac{R}{2} = R$.
\end{enumerate}
}
\end{bt}


% ==============================================================================
% PHẦN 2: ĐỀ THAM KHẢO TUYỂN SINH 10 QUẬN GÒ VẤP – ĐỀ 2 (7 BÀI)
% ==============================================================================

\begin{bt}%[9D6H18-2]
% nguon: Đề tham khảo Tuyển sinh 10 năm học 2025-2026 Sở GD&ĐT TP.HCM – Quận Gò Vấp - Đề 2 (Bài 1)
Cho parabol $(P)\colon y = 2x^2$.
\begin{enumerate}[a)]
    \item Vẽ đồ thị $(P)$ trên mặt phẳng tọa độ $Oxy$.
    \item Tìm tọa độ điểm $A$ thuộc đồ thị $(P)$ biết điểm $A$ có hoành độ bằng $-0{,}5$.
\end{enumerate}
\loigiai{
\begin{enumerate}[a)]
    \item Bảng giá trị tương ứng của hàm số $y = 2x^2$:
    \begin{center}
    \begin{tabular}{|c|c|c|c|c|c|}
        \hline
        $x$ & $-2$ & $-1$ & $0$ & $1$ & $2$ \\
        \hline
        $y = 2x^2$ & $8$ & $2$ & $0$ & $2$ & $8$ \\
        \hline
    \end{tabular}
    \end{center}
    Đồ thị $(P)$ là một parabol có đỉnh $O(0; 0)$, đối xứng qua trục tung $Oy$ và bề lõm quay lên trên.
    \item Điểm $A(x; y)$ thuộc $(P)$ có hoành độ $x = -0{,}5$.
    Thay $x = -0{,}5$ vào phương trình parabol, ta được:
    $y = 2 \cdot (-0{,}5)^2 = 2 \cdot 0{,}25 = 0{,}5$.
    Vậy tọa độ điểm cần tìm là $A(-0{,}5; 0{,}5)$.
\end{enumerate}
}
\end{bt}

\begin{bt}%[9D6H20-2]
% nguon: Đề tham khảo Tuyển sinh 10 năm học 2025-2026 Sở GD&ĐT TP.HCM – Quận Gò Vấp - Đề 2 (Bài 2)
Cho phương trình bậc hai $-3x^2 + 2x + 4 = 0$.
\begin{enumerate}[a)]
    \item Chứng minh phương trình luôn có hai nghiệm phân biệt $x_1, x_2$.
    \item Không giải phương trình, hãy tính giá trị của biểu thức $A = (x_1 - 2x_2)^2 - 3x_2^2$.
\end{enumerate}
\loigiai{
\begin{enumerate}[a)]
    \item Phương trình $-3x^2 + 2x + 4 = 0$ có các hệ số $a = -3$, $c = 4$.
    Vì $a \cdot c = (-3) \cdot 4 = -12 < 0$ nên phương trình luôn có hai nghiệm phân biệt trái dấu $x_1, x_2$.
    \item Áp dụng hệ thức Viète cho phương trình, ta có:
    $\begin{cases} S = x_1 + x_2 = -\dfrac{b}{a} = \dfrac{2}{3} \\ P = x_1 x_2 = \dfrac{c}{a} = -\dfrac{4}{3}. \end{cases}$
    Biến đổi biểu thức $A$:
    $A = (x_1 - 2x_2)^2 - 3x_2^2 = x_1^2 - 4x_1 x_2 + 4x_2^2 - 3x_2^2 = x_1^2 + x_2^2 - 4x_1 x_2 = (x_1 + x_2)^2 - 6x_1 x_2 = S^2 - 6P$.
    Thay các giá trị $S = \dfrac{2}{3}$ và $P = -\dfrac{4}{3}$ vào, ta được:
    $A = \left(\dfrac{2}{3}\right)^2 - 6 \cdot \left(-\dfrac{4}{3}\right) = \dfrac{4}{9} + 8 = \dfrac{76}{9}$.
\end{enumerate}
}
\end{bt}

\begin{bt}%[9D8H26-3]
% nguon: Đề tham khảo Tuyển sinh 10 năm học 2025-2026 Sở GD&ĐT TP.HCM – Quận Gò Vấp - Đề 2 (Bài 3)
Một chiếc hộp kín chứa $40$ quả bóng có cùng kích thước và khối lượng, gồm $5$ quả bóng màu đỏ, $12$ quả bóng màu xanh và các quả bóng còn lại màu vàng. Mỗi quả bóng được ghi một số tự nhiên phân biệt từ $1$ đến $40$.
\begin{enumerate}[a)]
    \item Lấy ngẫu nhiên $1$ quả bóng từ hộp. Tính xác suất của biến cố $A$: ``Số xuất hiện trên quả bóng lấy ra lớn hơn $30$''.
    \item Lấy ngẫu nhiên đồng thời $3$ quả bóng từ hộp. Tính xác suất của biến cố $B$: ``Lấy được $3$ quả bóng có màu khác nhau đôi một''.
\end{enumerate}
\loigiai{
Số quả bóng màu vàng có trong hộp là:
$40 - 5 - 12 = 23$ (quả bóng).
\begin{enumerate}[a)]
    \item Khi lấy ngẫu nhiên $1$ quả bóng, có $n(\Omega_1) = 40$ kết quả có thể xảy ra.
    Các số lớn hơn $30$ là từ $31$ đến $40$ ($10$ quả bóng), do đó $n(A) = 10$.
    Xác suất của biến cố $A$ là:
    $P(A) = \dfrac{n(A)}{n(\Omega_1)} = \dfrac{10}{40} = \dfrac{1}{4} = 0{,}25$.
    \item Khi lấy ngẫu nhiên $3$ quả bóng trong tổng số $40$ quả bóng, số kết quả có thể xảy ra là:
    $n(\Omega_2) = \mathrm{C}_{40}^3 = \dfrac{40 \cdot 39 \cdot 38}{6} = 9880$.
    Để lấy được $3$ quả bóng có màu khác nhau đôi một, ta phải lấy đúng $1$ quả bóng đỏ, $1$ quả bóng xanh và $1$ quả bóng vàng.
    Số cách chọn thuận lợi cho biến cố $B$ là:
    $n(B) = 5 \cdot 12 \cdot 23 = 1380$.
    Xác suất của biến cố $B$ là:
    $P(B) = \dfrac{n(B)}{n(\Omega_2)} = \dfrac{1380}{9880} = \dfrac{69}{494} \approx 0{,}1397$.
\end{enumerate}
}
\end{bt}

\begin{bt}%[9D1H3-4]
% nguon: Đề tham khảo Tuyển sinh 10 năm học 2025-2026 Sở GD&ĐT TP.HCM – Quận Gò Vấp - Đề 2 (Bài 4)
Một tổng kho sản xuất đồ gia dụng hiện có sẵn $900$ thùng hàng. Mỗi ngày phân xưởng xuất kho đúng $30$ thùng hàng để vận chuyển đi phân phối cho các đại lý bán lẻ.
\begin{enumerate}[a)]
    \item Gọi $T$ là số thùng hàng còn lại trong kho sau $c$ ngày xuất hàng. Hãy thiết lập công thức hàm số biểu diễn $T$ theo $c$.
    \item Biết giá trị mỗi thùng hàng là $2\,000\,000$ đồng và chi phí thuê xe tải vận chuyển mỗi chuyến chở $30$ thùng hàng mỗi ngày là $2\,500\,000$ đồng. Hỏi sau khi xuất bán hết toàn bộ số thùng hàng trong kho, doanh nghiệp thu về số tiền lãi ròng là bao nhiêu sau khi đã trừ toàn bộ chi phí vận chuyển?
\end{enumerate}
\loigiai{
\begin{enumerate}[a)]
    \item Sau $c$ ngày xuất hàng, số thùng hàng đã xuất kho là $30c$ (thùng).
    Công thức hàm số biểu diễn số thùng hàng còn lại trong kho là:
    $T(c) = 900 - 30c$ (thùng), với điều kiện $0 \le c \le 30$.
    \item Số ngày (tương ứng với số chuyến xe) cần thiết để phân phối hết $900$ thùng hàng là:
    $900 : 30 = 30$ (ngày/chuyến).
    Tổng số tiền thu được từ việc bán $900$ thùng hàng là:
    $900 \cdot 2\,000\,000 = 1\,800\,000\,000$ (đồng).
    Tổng chi phí vận chuyển trong $30$ chuyến là:
    $30 \cdot 2\,500\,000 = 75\,000\,000$ (đồng).
    Số tiền thực tế xưởng thu được sau khi trừ chi phí vận chuyển là:
    $1\,800\,000\,000 - 75\,000\,000 = 1\,725\,000\,000$ (đồng).
\end{enumerate}
}
\end{bt}

\begin{bt}%[9D1H3-4]
% nguon: Đề tham khảo Tuyển sinh 10 năm học 2025-2026 Sở GD&ĐT TP.HCM – Quận Gò Vấp - Đề 2 (Bài 5)
Thầy giáo biên soạn một đề thi kiểm tra môn Toán theo thang điểm $10$, gồm hai dạng câu hỏi: câu hỏi trắc nghiệm Đúng/Sai và câu hỏi trắc nghiệm nhiều lựa chọn, mỗi câu hỏi đều có giá trị điểm là $0{,}25$ điểm. Trong đề, số câu trắc nghiệm Đúng/Sai gấp ba lần số câu trắc nghiệm nhiều lựa chọn. Biết học sinh giải một câu trắc nghiệm Đúng/Sai mất $1$ phút và giải một câu nhiều lựa chọn mất $1{,}5$ phút. Hỏi học sinh có đủ thời gian làm bài trong tiết kiểm tra kéo dài $45$ phút hay không?
\loigiai{
Tổng số câu hỏi trong đề kiểm tra thang điểm $10$ là:
$10 : 0{,}25 = 40$ (câu).
Gọi số câu hỏi trắc nghiệm Đúng/Sai là $x$ và số câu hỏi nhiều lựa chọn là $y$ ($x, y \in \mathbb{N}^*$).
Theo đề bài, ta có hệ phương trình:
$\begin{cases} x + y = 40 \\ x = 3y \end{cases} \iff \begin{cases} 4y = 40 \\ x = 3y \end{cases} \iff \begin{cases} y = 10 \\ x = 30. \end{cases}$
Đề kiểm tra có $30$ câu Đúng/Sai và $10$ câu nhiều lựa chọn.
Thời gian học sinh cần dùng để hoàn thành toàn bộ bài kiểm tra là:
$T = 30 \cdot 1 + 10 \cdot 1{,}5 = 30 + 15 = 45$ (phút).
Thời gian này vừa đúng bằng thời lượng $45$ phút của tiết kiểm tra.
Vậy học sinh hoàn toàn đủ thời gian để hoàn thành bài kiểm tra.
}
\end{bt}

\begin{bt}%[9H0V31-3]
% nguon: Đề tham khảo Tuyển sinh 10 năm học 2025-2026 Sở GD&ĐT TP.HCM – Quận Gò Vấp - Đề 2 (Bài 6)
Một chai đựng nước khoáng có phần thân dưới dạng hình trụ với đường kính đáy bằng $6\text{ cm}$. Bạn An đo chiều cao cột nước trong chai khi đặt đứng được $9\text{ cm}$, sau đó nút chặt chai rồi lật ngược xuống thì phần hình trụ không chứa nước có chiều cao đo được là $7\text{ cm}$. Lấy $\pi \approx 3{,}14$.
\begin{enumerate}[a)]
    \item Tính thể tích lượng nước hiện có trong chai (kết quả làm tròn đến mililít).
    \item Tính tổng dung tích toàn bộ của chiếc chai (kết quả làm tròn đến mililít).
\end{enumerate}
\loigiai{
Bán kính đáy của phần chai hình trụ là:
$R = \dfrac{6}{2} = 3\text{ cm}$.
\begin{enumerate}[a)]
    \item Thể tích lượng nước trong chai là:
    $V_{\text{nước}} = \pi R^2 h_1 \approx 3{,}14 \cdot 3^2 \cdot 9 = 3{,}14 \cdot 81 = 254{,}34\text{ cm}^3 \approx 254\text{ ml}$.
    \item Khi lật ngược chai, phần không gian trống không chứa nước dồn về đáy hình trụ có chiều cao $h_2 = 7\text{ cm}$.
    Thể tích phần không gian trống trong chai là:
    $V_{\text{trống}} = \pi R^2 h_2 \approx 3{,}14 \cdot 3^2 \cdot 7 = 3{,}14 \cdot 63 = 197{,}82\text{ cm}^3 \approx 198\text{ ml}$.
    Tổng thể tích toàn bộ của chai nước là:
    $V_{\text{chai}} = V_{\text{nước}} + V_{\text{trống}} = 254{,}34 + 197{,}82 = 452{,}16\text{ cm}^3 \approx 452\text{ ml}$.
\end{enumerate}
}
\end{bt}

\begin{bt}%[9H9C29-3]
% nguon: Đề tham khảo Tuyển sinh 10 năm học 2025-2026 Sở GD&ĐT TP.HCM – Quận Gò Vấp - Đề 2 (Bài 7)
Từ điểm $A$ nằm bên ngoài đường tròn $(O; R)$, vẽ hai tiếp tuyến $AB, AC$ với đường tròn $(O)$ ($B, C$ là các tiếp điểm). Gọi $H$ là giao điểm của đoạn thẳng $OA$ và dây cung $BC$. Vẽ đường kính $BK$ của $(O)$, đoạn thẳng $AK$ cắt đường tròn $(O)$ tại điểm thứ hai là $E$.
\begin{enumerate}[a)]
    \item Chứng minh $OA \perp BC$ tại $H$ và $AH \cdot AO = AE \cdot AK$.
    \item Gọi $I$ là trung điểm của đoạn thẳng $AH$. Chứng minh $HE \perp EI$.
    \item Cho biết khoảng cách $OA = 2R$. Hãy tính diện tích tam giác $ABC$ theo $R$.
\end{enumerate}
\loigiai{
\begin{enumerate}[a)]
    \item Vì $AB, AC$ là hai tiếp tuyến cắt nhau nên $AB = AC$ và $OA$ là đường trung trực của đoạn thẳng $BC$, suy ra $OA \perp BC$ tại $H$.
    Trong tam giác vuông $ABO$ tại $B$ có đường cao $BH$, theo hệ thức lượng ta có:
    $AB^2 = AH \cdot AO$.
    Lại có $BK$ là đường kính nên $\widehat{BEK} = 90^\circ \implies BE \perp AK$ tại $E$.
    Trong tam giác vuông $ABK$ tại $B$ có đường cao $BE$, ta có:
    $AB^2 = AE \cdot AK$.
    Từ đó suy ra $AH \cdot AO = AE \cdot AK$ (đpcm).
    \item Vì $AH \cdot AO = AE \cdot AK \implies \dfrac{AH}{AK} = \dfrac{AE}{AO}$, kết hợp góc $\widehat{OAK}$ chung, suy ra $\triangle AHE \backsim \triangle AKO$ (c.g.c).
    Do đó $\widehat{AHE} = \widehat{AKO}$.
    Lại có tam giác $AHE$ vuông tại $E$ (hoặc xét vị trí trung điểm $I$ của cạnh huyền $AH$ trong tam giác vuông $AEH$ có trung tuyến $EI = IH = IA$), suy ra $\triangle IEH$ cân tại $I$, dẫn đến các góc phụ nhau và kết luận $HE \perp EI$.
    \item Trong tam giác vuông $ABO$ vuông tại $B$ có $OB = R$ và $OA = 2R$:
    $AB = \sqrt{OA^2 - OB^2} = \sqrt{(2R)^2 - R^2} = R\sqrt{3}$.
    Áp dụng hệ thức lượng trong tam giác vuông $ABO$:
    $AH = \dfrac{AB^2}{OA} = \dfrac{(R\sqrt{3})^2}{2R} = \dfrac{3R}{2}$.
    $BH = \dfrac{AB \cdot OB}{OA} = \dfrac{R\sqrt{3} \cdot R}{2R} = \dfrac{R\sqrt{3}}{2} \implies BC = 2BH = R\sqrt{3}$.
    Diện tích tam giác $ABC$ là:
    $S_{\triangle ABC} = \dfrac{1}{2} BC \cdot AH = \dfrac{1}{2} \cdot R\sqrt{3} \cdot \dfrac{3R}{2} = \dfrac{3R^2\sqrt{3}}{4}$.
\end{enumerate}
}
\end{bt}
